\chapter{Dataset Preparation}
\label{ch:dataset}

\section{Dataset Preparation}

\textbf{Data Source:} The audio recordings used for training the proposed model were collected from two primary sources: publicly available recordings from YouTube and recordings performed by the authors. The YouTube recordings were selected to obtain high-quality, low-noise samples, while the self-recorded flute performances were included to introduce realistic environmental noise and recording variations. Combining these two sources resulted in a more diverse dataset, enabling the model to learn from both clean and noisy audio conditions and thereby improving its robustness.

After collection, each audio clip was manually categorized according to its corresponding swara (Sa, Re, Ga, Ma, Pa, Dha, and Ni). The recordings were then organized into separate folders, with each folder containing audio samples of a single swara. This folder-based organization simplified data management and facilitated the preprocessing and training stages of the proposed framework. The directory structure of the dataset is illustrated in Figure~\ref{fig:folder-structure}.

\figplaceholder{folder_structure.png}{Folder structure of dataset}{fig:folder-structure}

Each swara class contains more than 100 audio recordings, providing a sufficiently balanced dataset for model training. Furthermore, the recordings exhibit a range of note durations, allowing the dataset to capture natural temporal variations in flute performances. The class distribution and the duration distribution of the individual swara recordings are presented in Figure~\ref{fig:class-dist} (originally presented as two side-by-side subplots in the source document).

\figplaceholder{class distribution.png}{Class Distribution}{fig:class-dist}
\figplaceholder{duration.png}{Duration Distribution of Collected Swaras}{fig:duration}

\begin{table}[H]
\centering
\caption{Individual Swara Dataset}
\label{tab:individual-swara}
\begin{tabular}{ll}
\toprule
\textbf{Metric} & \textbf{Value} \\
\midrule
Total Recordings & 748 \\
Average Recording per Class & 106.86 \\
Number of Swara Class & 7 \\
Average Duration & 1.137 (sec) \\
Median Duration & 0.968 (sec) \\
Minimum Duration & 0.305 (sec) \\
Maximum Duration & 13.62 \\
Standard Deviation of Duration & 1.014 (sec) \\
\bottomrule
\end{tabular}
\end{table}

\subsection*{Synthetic Melody Dataset Creation}

To facilitate the evaluation of the proposed swara segmentation and classification framework, a synthetic melody dataset was generated by sequentially combining isolated flute recordings of the seven fundamental swaras (Sa, Re, Ga, Ma, Pa, Dha, and Ni). Each melody was constructed as an ordered sequence of individual note recordings, with every note assigned a unique sequence index to preserve its temporal position within the melody. During the generation process, the onset and offset timestamps of each swara were recorded, producing frame-level ground-truth annotations that precisely describe the temporal boundaries of every note in the synthesized melody. These annotations were stored in a metadata file containing the melody filename, note sequence index, swara label, start time, end time, transition styles and DSP applied. A portion of the melody dataset is presented in Table~\ref{tab:synthetic-melody}.

\begin{table}[H]
\centering
\caption{Synthetic melody dataset (sample)}
\label{tab:synthetic-melody}
\resizebox{\textwidth}{!}{%
\begin{tabular}{lllllll}
\toprule
\textbf{Melody file} & \textbf{Note\_sequence} & \textbf{Swara\_label} & \textbf{Start\_time} & \textbf{End\_time} & \textbf{Transition\_style} & \textbf{Applied\_dsp} \\
\midrule
flute\_melody\_0.wav & 0 & Sa & 0 & 0.4043 & phrase\_start & speed(1.05)+vibrato+gain(1.02) \\
flute\_melody\_0.wav & 1 & Ma & 0.3325 & 0.9921 & slurred\_crossfade & pitch(0.17)+speed(1.16)+vibrato+gain(0.63)+noise \\
flute\_melody\_0.wav & 2 & Pa & 1.2282 & 3.4976 & articulated\_with\_breath & pitch(1.06)+speed(0.82)+gain(1.23)+noise \\
flute\_melody\_0.wav & 3 & Ni & 3.6154 & 4.678 & articulated\_with\_breath & pitch(1.38)+speed(1.08)+gain(0.82)+noise \\
flute\_melody\_0.wav & 4 & Ni & 4.9613 & 6.0312 & articulated\_with\_breath & speed(0.99)+gain(1.22) \\
\bottomrule
\end{tabular}%
}
\end{table}

The synthetic melodies were generated by randomly selecting note recordings from the pool of isolated swara samples. For each melody, the sequence of swaras, the number of notes, and the specific recording chosen for each swara were determined programmatically to maximize diversity while ensuring that all seven swaras were represented throughout the dataset. Since multiple recordings were available for each swara, the generator selected different source recordings across melodies, preventing excessive repetition of identical note samples. Furthermore, the duration of each note within a melody was varied by assigning different temporal lengths, resulting in melodies with diverse rhythmic structures and overall durations. This randomization strategy produced a wide range of melodic patterns while preserving accurate ground-truth annotations for every note. The diverseness of notes per melody and duration is shown in Figures~\ref{fig:melody-duration} and~\ref{fig:notes-per-melody}.

\figplaceholder{fig4_melody_duration_distribution.png}{Melody Duration Distribution}{fig:melody-duration}
\figplaceholder{fig5_notes_per_melody.png}{Distribution of Notes per Synthetic Melody}{fig:notes-per-melody}

To further improve the realism of the synthesized melodies, multiple transition techniques and digital signal processing (DSP) augmentations were incorporated during the generation process. The metadata indicates that transition styles such as direct transitions, crossfades, and silence-separated note boundaries were applied between consecutive swaras to simulate different articulation patterns observed in flute performances. In addition, several DSP-based augmentations, including pitch shifting, time stretching, gain adjustment, additive noise, vibrato, and reverberation, were selectively applied to individual notes or melody segments. The specific augmentation applied to each note was recorded in the metadata, allowing complete traceability of the synthesis process. These variations increase the acoustic diversity of the dataset and improve its suitability for evaluating the robustness and generalization capability of the proposed swara segmentation and classification framework. Figure~\ref{fig:dsp-usage} visualizes the usage of different DSP augmentations.

\figplaceholder{fig6_dsp_augmentation_usage.png}{Count of different DSP augmentations used while creating the Synthetic Melody Dataset}{fig:dsp-usage}

The resulting metadata file serves as a comprehensive annotation source for the synthetic dataset, allowing direct computation of melody-level statistics, note durations, swara occurrence frequencies, transition patterns, and temporal distributions. These detailed annotations provide reliable ground truth for quantitative evaluation of both segmentation accuracy and swara classification performance.

The synthetic melody dataset was developed to address one of the primary challenges in automatic swara segmentation: the scarcity of accurately annotated continuous musical recordings. Creating a large-scale annotated dataset of flute performances requires substantial time and expert effort, as every note boundary and swara label must be manually identified and verified. Such a process is not only labor-intensive but also difficult to scale for different instruments and musical styles. To overcome this limitation, a synthetic melody generation pipeline was designed to automatically produce continuous melodies with precise ground-truth annotations, eliminating the need for extensive manual labeling.

Although the proposed framework is evaluated using Indian classical flute recordings, the underlying dataset generation pipeline is designed to be instrument-independent and can be adapted to any monophonic instrument for which isolated note recordings are available. By automatically generating continuous melodies from individual note samples and simultaneously producing accurate temporal annotations, the proposed approach significantly reduces the effort required to prepare training data. This enables the development of robust swara segmentation and classification models without relying on hundreds of hours of manually annotated recordings, making the framework more scalable and practical for future research on monophonic musical instruments.
